\addcontentsline{toc}{chapter}{Abstract}
% A short summary of the project, the report and what the reader
% will get from reading it.


\noindent
Efficient heat production and energy management are becoming increasingly important in modern district heating systems, particularly in situations where operational costs, electricity prices, and energy demand continuously fluctuate. The project focused on developing a software solution capable of optimizing heat production schedules while ensuring stable heat delivery to consumers and improving overall economic performance. The system analyses production unit configurations, heat demand data, and electricity market conditions to determine the most cost-effective production strategy across different operational scenarios.\par
The work was carried out using Scrum methodology and iterative software development practices. Multiple sprints were used to organise tasks, manage collaboration, and continuously improve the system through planning sessions, reviews, and retrospectives. Alongside the project management aspects, significant attention was given to software architecture, system modelling, implementation, testing, and code quality.\par  
The document presents the complete development process, beginning with the project background, objectives, and requirements, followed by the Scrum implementation and sprint activities. It also explains the technical structure of the solution through UML diagrams, architectural descriptions, frontend and backend implementation details, and data management strategies. In addition, reflections on teamwork, testing, refactoring, code reviews, and requirements engineering are discussed to evaluate both the technical and collaborative outcomes of the project.\par
By exploring the different stages of development and the decisions made throughout the process, the reader gains insight into how Agile methodologies and software engineering practices can be applied to create a functional optimization system for district heating production.
